\begin{abstract}

We show that a variety of modern deep learning tasks exhibit
a ``double-descent'' phenomenon where, as we increase model size, performance first gets \textit{worse} and then gets better.
Moreover, we show that double descent occurs not just as a function of model size, but also as a function of the number of training epochs.
We unify the above phenomena by defining
a new complexity measure we call the
\emph{effective model complexity}
and conjecture a generalized double descent with respect to this measure.
%which appears to determine
%the location of this phase transition, 
Furthermore,
our notion of model complexity
allows us to identify
certain regimes where increasing (even quadrupling) the number of
train samples actually \emph{hurts} test performance. 

\iffalse
We show that a variety of natural learning tasks exhibit a "double-descent" phenomena, where, as we increase model size, performance first gets \textit{worse} and then gets better.

%We also observe the double descent phenomenon as we increase the size of the training set.  

%We present experimental findings in modern deep learning
%that both challenge and unify conventional wisdoms on
%the relation of test error to model complexity.
%We show that a variety of natural learning tasks exhibit a
%\emph{phase transition phenomenon} between ``classical'' and ``modern'' regimes.
%dependence of performance on complexity.

In one parameter regime, test error follows the  classical ``bias-variance curve'', where bigger models can hurt generalization,
while in another regime bigger models only help.
We define a new complexity measure we call the \emph{effective model complexity} which appears to determine the location of this phase transition, 
by incorporating the optimization algorithm into traditional notions of model complexity.
%and unifies traditional notions of model complexity with the length of training. 

As a consequence we establish the existence of the ``double descent'' phenomenon, introduced by \cite{belkin2018reconciling}, in a variety of modern datasets and architectures, 
including CIFAR 10/100 with ResNet18s/CNNs, and IWSLT'14/WMT'14 with Transformers.
We show that double descent occurs not just as a function of model size, but also as a function of the number of training epochs.
Furthermore, the relation between effective model complexity and  the number of samples allows us to identify regimes where increasing (even quadrupling!) the number of 
samples only \emph{hurts} test performance. 

\fi

% Before Ilya comments
\iffalse
We present experimental findings in modern deep learning
that both challenge and unify conventional wisdoms on
the relation of test error to model complexity.
We show that a variety of natural learning tasks and training procedures exhibit a
\emph{phase transition phenomenon} between ``classical'' and ``modern'' regimes.
%dependence of performance on complexity.
In one parameter regime, test error follows the  classical ``bias-variance curve'', where bigger models can hurt generalization,
while in another regime bigger models only help.
We define a new complexity measure we call the \emph{effective model complexity} which appears to determine the location of this phase transition, 
by incorporating the optimization algorithm into traditional notions of model complexity.
%and unifies traditional notions of model complexity with the length of training. 

As a consequence we establish the existence of the ``double descent'' phenomenon, introduced by \cite{belkin2018reconciling}, in a variety of modern datasets and architectures, 
including CIFAR 10/100 with ResNet18s/CNNs, and IWSLT'14/WMT'14 with Transformers.
We show that double descent occurs not just as a function of model size, but also as a function of the number of training epochs.
Furthermore, the relation between effective model complexity and  the number of samples allows us to identify regimes where increasing (even quadrupling!) the number of 
samples only \emph{hurts} test performance. 
\fi



\iffalse
OLD abstract:

We empirically demonstrate a ``double descent'' phenomenon for state of the art neural networks for an array of standard tasks.
Specifically, we show for a variety of models, training procedures and datasets, the following occurs---with increase of either model size or train time, the classifier undergoes the following transition:
\begin{enumerate}
    \item Initially, as model size or train time increases, the test error \emph{decreases}.
    \item At some point, as model size or train time continue to grow, the test error \emph{increases} reaching a peak around the ``interpolation threshold'' where the model and training procedure are just barely able to achieve $\approx 0$ train error.
    \item As we continue to increase model size or train time beyond this point, the test error undergoes a \emph{second descent}, decreasing with train time or model size.
\end{enumerate}
 
Such a ``double descent'' for varying model size was recently
%demonstrated by 
proposed by
\cite{belkin2018reconciling}.
%for random Fourier Features and two-layer neural networks on MNIST.
We show that this extends to many modern settings,
and applies to both varying model-size to varying train-time.
%extends to many modern settings, and generalized from model-size
%to 
Moreover, we introduce a new notion of ``Effective Model Complexity''
that unifies both these double descent phenomenon. % in both model size and train time.
%and we conjecture a generalized double-descent with respect to this complexity.
We also show that double descent can give rise to \emph{non-monotonicity in sample size}: there are regimes where increasing the number of training samples can lead to \emph{worse} test performance.
\fi
\end{abstract}